% See https://holatex.app/playground.html# for an online viewer.

\documentclass[tikz]{standalone} % Standalone: Only graphics export.
\usepackage[european,siunitx]{circuitikz}
\usetikzlibrary{arrows.meta}
\usepackage{listings}
\lstset{basicstyle=\footnotesize\ttfamily}
\begin{document}

  \newcommand{\eq}{=}
  \tikzset{ameter/.style={rmeterwa, t=A}}
  \begin{circuitikz}
  %\begin{circuitikz}[american voltages, american currents]

    \tikzset{>=Triangle} % Define tip of arrows drawn with [->]
    \ctikzset{voltage=straight}

    \def\sizeY{\lR}         % Heights Uin
    \def\lR{2.5}            % Length of device like R
    \def\dxOp12{0.3*\lR}    % Delta x between both OP input lines.
    \def\xOp1{\lR}          % X-coordinate of left OP input line.
    
    \def\lUio{0.8}          % Length of I/O connectors
    \def\gapUio{0.3}        % Gap between I/O connector and voltage arrow

    \coordinate(origin) at (0,0);
    
    % Isolated input node.
    \draw (0,\sizeY) coordinate(in)
      to[short,o-] (in);

%      
%      to[R,l=$R$,o-*] ++(\lR,0) coordinate(BotLftFb)
%      -- ++(0,\dFbY) to[short,-*] ++(\dRcX,0) coordinate(lftRc1)
%      -- ++(0,-\dRcY)
%      to[C,l=$C_{t_1}$,v=$U_{C_{t_1}}$] ++(\lR,0)
%      to[ameter,i=$i_{C_{t_1}}$] ++(\lR,0)
%      to[short,-*] ++(0,\dRcY) coordinate(rhtRc1)
%      to[short,-*] ++(\dRcX,0) coordinate(lftRc2) -- ++(0,-\dRcY)
%      to[C,l=$C_{t_2}$,v=$U_{C_{t_2}}$] ++(\lR,0)
%      to[ameter,i=$i_{C_{t_2}}$] ++(\lR,0)
%      to[short,-*] ++(0,\dRcY) coordinate(rhtRc2)
%      to[R,l=$R_u$] ++(\lR,0) coordinate(rhtRu);
%    
%    % Upper left part of feedback line.
%    \draw (lftRc1) -- ++(0,\dRcY)
%    to[R,l=$R_{t_1}$] ++(\lR,0) coordinate(rhtRt1)
%    -- (rhtRt1 -| rhtRc1) -- (rhtRc1);
%    
%    % Upper right part of feedback line.
%    \draw (lftRc2) -- ++(0,\dRcY)
%    to[R,l=$R_{t_2}$] ++(\lR,0) coordinate(rhtRt2)
%    -- (rhtRt2-|rhtRc2) -- (rhtRc2);
    
    % OP
    \draw (\xOp1,0)
      to[short,*-] ++(0,\sizeY)
      -- ++(\dxOp12,0)
      node[plain amp,anchor=-](OP){$OP$}
      (OP.out)
      to[short,-o] ++(\lUio,0) coordinate(out)
      (OP.+)
      to[short,-*] (origin-|OP.+)
      node[ground]{};
    
    % Ground line.
    \draw (origin) to[short,o-o] (origin-|out) coordinate(botOut);
%      to[short,-o] ($(origin-|BotRhtFb)+(\lUio,0)$) coordinate(gndOut);
%    
%    % OP to ground symbol.
%    \draw (OP.+) -- (origin-|OP.-) node[ground]{};
%
    % Voltage arrow for input voltage U1.
    \draw[->] ($(in)+(0,-\gapUio)$) -- node[midway,left]{$U_1$} ($(origin)+(0,\gapUio)$);

%    % Connectors for output voltage.
%    \draw (BotRhtFb) to[short, *-o] ++(\lUio,0) coordinate(out);

    % Voltage arrow for output voltage.
    \draw[->] ($(out)+(0,-\gapUio)$) -- node[midway,right]{$U_2$}
      ($(botOut)+(0,\gapUio)$);

  \end{circuitikz}

\end{document}
